%%________________________________________________________________________
%% LEIM | PROJETO -- Back To School
%% Apendices A e B
%%________________________________________________________________________




%%________________________________________________________________________
\chapter{Gestão de Versões}
\label{ap:gestaoVersoes}
%%________________________________________________________________________

Esta secção descreve o processo utilizado na gestão de versões do projeto ao longo do seu desenvolvimento, incluindo a arquitetura de repositórios utilizada, a estratégia de ramificação adotada e o fluxo de sincronização entre os elementos do grupo.

Para o controlo de versões do projeto foi utilizado o sistema Git, com o repositório remoto alojado na plataforma GitHub. O uso do Git revelou-se essencial, uma vez que permitiu que cada elemento do grupo mantivesse um clone local do repositório, sincronizado com o repositório remoto através de operações de \textit{push} e \textit{pull}. Adicionalmente, uma vez que o projeto consistiu no desenvolvimento de um jogo em Unity, foi também utilizada a extensão Git LFS, permitindo que os ficheiros de maiores dimensões, como por exemplo \textit{assets} referentes a modelos 3D, fossem geridos separadamente do histórico principal do repositório.

Relativamente à estratégia de ramificação, optou-se por uma abordagem simplificada e direta, tendo o projeto sido desenvolvido numa única \textit{branch} principal (\textit{main}). Este foi o único ramo definido e ativamente utilizado ao longo de todo o desenvolvimento, com o objetivo de manter, de forma contínua, a versão mais recente do trabalho realizado.




%%________________________________________________________________________
\chapter{Apoio de Inteligência Artificial ao Desenvolvimento}
\label{ap:apoioIA}
%%________________________________________________________________________

Ao longo do desenvolvimento do projeto, foram utilizadas ferramentas de inteligência artificial generativa como apoio ao trabalho. Numa primeira fase, a IA serviu para alinhar as ideias do grupo sobre os pontos a implementar prioritariamente, ajudando a organizar e a estruturar o planeamento inicial do sistema de jogo. Posteriormente, foi ainda utilizada para apoiar a compreensão de conceitos técnicos e tecnologias envolvidas no projeto, bem como no processo de \textit{debugging}, na identificação de possíveis causas de erros, na sugestão de correções ao código e no apoio ao desenvolvimento de novas implementações.

Sempre que sugestões geradas por IA eram integradas no projeto, o grupo teve o cuidado de as testar e validar antes de as considerar definitivas. Esta verificação foi particularmente relevante em partes mais sensíveis do sistema, como a sincronização em rede e o comportamento dos NPCs, onde um erro não identificado poderia comprometer a experiência multijogador.

Em todos estes casos, a IA teve um papel de apoio ao processo de desenvolvimento, sem substituir as decisões técnicas, os testes e as validações realizados pelo grupo ao longo do trabalho.
